\documentclass[11pt]{article}

\usepackage{amsmath}
\usepackage{amssymb}
\usepackage{booktabs}
\usepackage[margin=1in]{geometry}
\usepackage[round,authoryear]{natbib}
\usepackage{xcolor}
\usepackage{hyperref}
\hypersetup{
  colorlinks=true,
  linkcolor=blue!50!black,
  citecolor=blue!50!black,
  urlcolor=blue!50!black,
}

% --- consistent macros ------------------------------------------------------
\newcommand{\T}{\mathsf{T}}                    % transpose
\newcommand{\dx}{\Delta x}                     % log-wavelength shift
\newcommand{\half}{\tfrac{1}{2}}
\newcommand{\rfft}{\mathrm{rfft}}
\newcommand{\irfft}{\mathrm{irfft}}
\newcommand{\diag}{\operatorname{diag}}
\newcommand{\argmin}{\operatorname*{arg\,min}}
\newcommand{\Var}{\operatorname{Var}}
\newcommand{\velscale}{\mathrm{velscale}}

\title{A Separable Framework for Joint Radial Velocity,\\
Template Weight, Continuum, and Spectral Resolution\\
Fitting of Stellar Spectra}
\author{}
\date{}

\begin{document}
\maketitle

\begin{abstract}
We derive the spectral-fitting framework implemented in this project, which
simultaneously recovers a radial velocity $v$, the weights $\mathbf{w}$ of a
learned (PCA) template basis, a smooth Legendre continuum $\mathbf{c}$, and
optionally the instrumental spectral resolution $R$, from a single observed
stellar spectrum. The key structural fact, exploited throughout, is
\emph{separability}: working in log-wavelength $x=\ln\lambda$ and in log-flux,
the radial velocity becomes a uniform translation, the resolution becomes a
fixed-width Gaussian, and the continuum becomes additive, so that at fixed values
of the nonlinear parameters $(v,R)$ the model is \emph{linear} in
$\boldsymbol\theta=(\mathbf{w},\mathbf{c})$ and admits an exact weighted
least-squares solution. We build the model in seven steps: the linear template
fit; the radial-velocity translation and per-shift solve (the cross-correlation
of \citealt*{TonryDavis1979}, generalised to a template basis); FFT acceleration
of the entire velocity scan; the convexity structure of the joint problem
(convex in the linear block, multimodal in $v$); the additive Legendre continuum
(in the spirit of the multiplicative/additive polynomials of pPXF,
\citealt*{Cappellari2017}); the on-the-fly Gaussian broadening for resolution,
together with an honest accounting of the log-space convolution approximation;
and finally the actual optimiser --- a coarse Fourier scan to locate the global
basin followed by a JAX-autodiff damped-Newton refinement of the profiled
$\chi^2$, with covariances from its Hessian. The PCA-template basis and the
external resolution/redshift scan parallel the classification approach of
\citet*{Bolton2012}.
\end{abstract}

% ============================================================================
\section*{Notation}
% ============================================================================
All symbols used in this document are collected in Table~\ref{tab:notation} and
are not redefined elsewhere. Vectors are bold lowercase (the log-flux data $\mathbf{Y}$ is the lone
bold-uppercase vector), matrices bold uppercase; the per-pixel index is $i=1,\dots,n$; a hat denotes an estimate.

\begin{table}[h]
\centering
\small
\begin{tabular}{@{}llp{0.55\textwidth}@{}}
\toprule
Symbol & Type & Meaning \\
\midrule
$\lambda$            & --        & wavelength \\
$c$                  & scalar    & speed of light \\
$x=\ln\lambda$       & --        & log-wavelength coordinate \\
$x_i$                & pixels    & uniform log-wavelength grid, $i=1,\dots,n$ \\
$\delta$             & scalar    & constant pixel spacing in $x$, $\;\delta=x_{i+1}-x_i$ \\
$n$                  & scalar    & number of pixels \\[2pt]
$v$                  & scalar    & radial velocity \\
$\dx=\ln(1+v/c)$     & scalar    & log-wavelength shift induced by $v$ \\
$p=\dx/\delta$       & pixels    & translation in pixels (the ``lag'') \\[2pt]
$d_i$                & vector    & observed (linear) flux in pixel $i$ \\
$\sigma_i$           & vector    & $1\sigma$ noise on $d_i$ \\
$Y_i=\ln d_i$        & vector    & log-flux (the modelled quantity) \\
$C_i=(\sigma_i/d_i)^2$ & vector  & variance of $Y_i$; weight $C_i^{-1}=0$ if masked \\
$\mathbf{C}=\diag(C_i)$ & matrix & diagonal data covariance of $\mathbf{Y}$; weight $\mathbf{C}^{-1}$ \\[2pt]
$T_0=\mu$            & vector    & mean rectified log-flux template \\
$T_k=\phi_k$         & vector    & $k$-th rectified-log-flux PCA component, $k=1,\dots,K$ \\
$\mathbf{T}$         & $n\times(K{+}1)$ & template design block, columns $[\mu,\phi_1,\dots,\phi_K]$ \\
$K$                  & scalar    & number of PCA components \\
$\mathbf{w}=(w_0,\dots,w_K)$ & vector & template weights ($w_0$ multiplies $\mu$) \\[2pt]
$P_m$                & vector    & $m$-th Legendre continuum column, $m=0,\dots,L$ \\
$\mathbf{P}$         & $n\times(L{+}1)$ & continuum design block \\
$L$                  & scalar    & continuum (Legendre) order \\
$\mathbf{c}=(c_0,\dots,c_L)$ & vector & continuum coefficients \\[2pt]
$\tau_i = -Y_i^+$    & vector    & non-negative optical depth, $\tau_i=\max(-Y_i,0)$, used in NMF training \\
$\boldsymbol\tau$    & $n_{\rm spec}\times n_{\rm pix}$ & optical-depth matrix of the training library \\
$\mathbf{W}_{\rm NMF}$ & $n_{\rm spec}\times K$ & NMF activation matrix (training, $\geq 0$) \\
$\mathbf{H}_{\rm NMF}$ & $K\times n_{\rm pix}$ & NMF component matrix (training, $\geq 0$) \\[2pt]
$S_j,\;\mathbf{S}$   & vector, $n\times(J{+}1)$ & telluric templates / block (observed frame, not RV-shifted) \\
$\mathbf{u}=(u_0,\dots,u_J)$ & vector & telluric weights \\
$p_{\rm tell}$       & pixels    & telluric lag (instrument wavelength zero-point) \\[2pt]
$\mathbf{A}$         & $n\times D$ & full design $[\,\mathbf{T}\mid\mathbf{P}\,]$ (or $[\,\mathbf{T}\mid\mathbf{S}\mid\mathbf{P}\,]$ with tellurics) \\
$\boldsymbol\theta=(\mathbf{w},\mathbf{c})$ & vector & full linear parameter vector \\
$\mathbf{N}=\mathbf{A}^\T\mathbf{C}^{-1}\mathbf{A}$ & matrix & normal matrix \\
$\mathbf{r}=\mathbf{A}^\T\mathbf{C}^{-1}\mathbf{Y}$ & vector & right-hand side \\
$\mathbf{M}_{TT},\mathbf{M}_{TP},\mathbf{M}_{PP}$ & matrix & blocks of $\mathbf{N}$ \\
$\mathbf{b}_T,\mathbf{b}_P$ & vector & blocks of $\mathbf{r}$ \\[2pt]
$R=\lambda/\Delta\lambda$ & scalar & spectral resolving power \\
$\sigma_v$           & scalar    & LSF velocity dispersion \\
$\sigma_{\rm pix}$   & scalar    & LSF width in pixels on the log grid \\
$\velscale=c\,\delta$ & scalar   & velocity per pixel \\
$f=2\sqrt{2\ln 2}$   & scalar    & FWHM-to-$\sigma$ conversion factor \\[2pt]
$\omega$             & vector    & Fourier frequency (cycles per pixel) \\
$\rfft,\irfft$       & --        & real forward / inverse FFT \\
$\chi^2$             & scalar    & weighted least-squares objective \\
$\widehat{(\cdot)}$  & --        & estimate of a quantity \\
\bottomrule
\end{tabular}
\caption{Strict notation used throughout. Symbols are defined once here.}
\label{tab:notation}
\end{table}

% ============================================================================
\section{Setup and the linear template fit}
\label{sec:setup}
% ============================================================================
We begin by resampling both the observed spectrum and the templates onto a single
uniform grid in log-wavelength, $x_i = x_1 + (i-1)\delta$ with $i=1,\dots,n$ and
constant spacing $\delta$. Two facts motivate this grid: a Doppler shift is a
uniform translation in $x$ (Section~\ref{sec:rv}), and at constant resolving
power the line-spread function (LSF) has a constant width in $x$
(Section~\ref{sec:resolution}).

We model the spectrum in \emph{log-flux}, $Y_i=\ln d_i$. This is the natural space
for two reasons. First, the observed flux is the product of a smooth continuum and
a normalised (rectified) spectrum, and the logarithm turns this product into a
sum, making the continuum an \emph{additive} term (Section~\ref{sec:continuum}).
Second, the template basis is built in log-rectified flux, as follows. Each model
spectrum is converted to log-flux, a robust low-order Legendre fit to the upper
(continuum) envelope is subtracted to \emph{rectify} it, and a PCA is performed on
the rectified spectra. This yields a mean rectified log-flux spectrum $\mu(x)$ and
orthonormal principal components $\phi_k(x)$, $k=1,\dots,K$, that capture the
\emph{line} information. We collect them as the columns of the template block,
\begin{equation}
  \mathbf{T} = \big[\,\mu \;\big|\; \phi_1 \;\big|\; \cdots \;\big|\; \phi_K\,\big],
  \qquad T_0\equiv\mu,\quad T_k\equiv\phi_k,
\end{equation}
so that $\mathbf{T}$ has $K+1$ columns. The mean $\mu$ is treated as a
free-weight template (column $T_0$, weight $w_0$), so it is Doppler-shifted like
any component.

The simplest fitting problem models the observed log-rectified flux as a linear
combination of the templates,
\begin{equation}
  Y(x) \;=\; \sum_{k=0}^{K} w_k\, T_k(x),
  \label{eq:linmodel}
\end{equation}
with weight vector $\mathbf{w}=(w_0,\dots,w_K)$.

\paragraph{Data covariance.}
The noise is naturally specified on the linear flux $d_i$ with dispersion
$\sigma_i$. Propagating to log-flux via $Y_i=\ln d_i$ and the first-order
(delta-method) linearisation gives the variance of $Y_i$,
\begin{equation}
  \sigma_{Y,i}\;\approx\;\left|\frac{\partial \ln d}{\partial d}\right|\sigma_i
  \;=\;\frac{\sigma_i}{d_i},
  \qquad\Longrightarrow\qquad
  C_i \;=\; \sigma_{Y,i}^2 \;=\; \left(\frac{\sigma_i}{d_i}\right)^2 ,
  \label{eq:weights}
\end{equation}
so the inverse-variance weight is $C_i^{-1}=(d_i/\sigma_i)^2$. Bad or out-of-range
pixels (non-finite $d_i$ or $\sigma_i$, $d_i\le 0$, masked pixels) are excised by
setting $C_i^{-1}=0$ (equivalently $C_i\to\infty$); we write $\mathbf{C}=\diag(C_i)$
and let $n_{\rm good}=\#\{i:C_i^{-1}>0\}$ be the number of usable pixels.

\paragraph{Weighted least squares.}
With the design matrix taken to be the template block $\mathbf{T}$, the objective
is the weighted residual sum of squares,
\begin{equation}
  \chi^2(\mathbf{w}) \;=\; (\mathbf{Y}-\mathbf{T}\mathbf{w})^\T \mathbf{C}^{-1}\,(\mathbf{Y}-\mathbf{T}\mathbf{w}).
  \label{eq:chi2lin}
\end{equation}
This is a quadratic in $\mathbf{w}$; setting $\nabla_{\mathbf w}\chi^2=0$ yields
the normal equations
\begin{equation}
  \mathbf{N}\,\mathbf{w} \;=\; \mathbf{r},
  \qquad
  \mathbf{N} \;=\; \mathbf{T}^\T\mathbf{C}^{-1}\mathbf{T},
  \qquad
  \mathbf{r} \;=\; \mathbf{T}^\T\mathbf{C}^{-1}\mathbf{Y},
  \label{eq:normaleq}
\end{equation}
with the unique solution and minimum value
\begin{equation}
  \widehat{\mathbf{w}} \;=\; \mathbf{N}^{-1}\mathbf{r},
  \qquad
  \chi^2(\widehat{\mathbf{w}}) \;=\; \mathbf{Y}^\T\mathbf{C}^{-1}\mathbf{Y}
  \;-\; \mathbf{r}^\T \mathbf{N}^{-1}\mathbf{r}.
  \label{eq:chi2profiled}
\end{equation}
The Hessian of \eqref{eq:chi2lin} is $2\mathbf{N}=2\,\mathbf{T}^\T\mathbf{C}^{-1}\mathbf{T}\succeq 0$,
so the problem is a convex quadratic; it is \emph{strictly} convex --- with a
unique minimiser --- if and only if $\mathbf{T}^\T\mathbf{C}^{-1}\mathbf{T}$ is
positive definite, which requires the weighted template columns to be linearly
independent. If they are not (a (near-)singular $\mathbf{N}$, e.g.\ from
collinear components or too few good pixels), the solve is regularised: a Tikhonov
ridge $\mathbf{N}\!\to\!\mathbf{N}+\rho\mathbf{I}$, a truncated SVD, or a
non-negativity constraint (NNLS) when physical positivity of the weights is
desired.

% ============================================================================
\section{Template basis: PCA versus NMF}
\label{sec:basis}
% ============================================================================
The template columns $T_k$ in \eqref{eq:linmodel} can be constructed from a
training library by two methods: PCA or NMF. Both operate on the
$n_{\rm spec}\times n_{\rm pix}$ rectified log-flux matrix $\mathbf{R}$, whose rows
are training spectra after continuum removal. At test time the same separable
machinery is used regardless of which basis was chosen; the difference lies entirely
in how the basis vectors are built and in the constraint placed on the weights at
the final solve.

\subsection*{PCA basis}

PCA finds the $K$-dimensional affine subspace of minimum weighted reconstruction
error. The sample mean $\mu$ is computed and subtracted; a truncated SVD of the
mean-centred matrix yields $K$ orthonormal components $\phi_k$ and their associated
eigenvalues $\lambda_k$ (one per component), ordered by decreasing explained
variance. At test time the weight vector $\mathbf{w}$ is solved by unconstrained
weighted least squares \eqref{eq:normaleq}; each $w_k$ is free to be positive or
negative.

\subsection*{NMF basis}

NMF exploits the sign structure of absorption spectra. Because rectified log-flux is
non-positive wherever stellar opacity is nonzero ($Y_i \leq 0$ in the lines), we
define the non-negative optical-depth matrix
\begin{equation}
  \tau_{si} \;=\; \max(-R_{si},\,0) \;\geq\; 0
\end{equation}
and factorise it as
\begin{equation}
  \boldsymbol\tau \;\approx\; \mathbf{W}_{\rm NMF}\,\mathbf{H}_{\rm NMF},
  \qquad
  \mathbf{W}_{\rm NMF} \in \mathbb{R}^{n_{\rm spec}\times K}_{\geq 0},
  \quad
  \mathbf{H}_{\rm NMF} \in \mathbb{R}^{K\times n_{\rm pix}}_{\geq 0},
  \label{eq:nmffac}
\end{equation}
minimising $\|\boldsymbol\tau - \mathbf{W}_{\rm NMF}\mathbf{H}_{\rm NMF}\|_F^2$
over non-negative factors by Lee--Seung multiplicative updates
\citep*{LeeSeung1999}, initialised with the non-negative double SVD (NNDSVD) of
\citet*{BoutsidisGallopoulos2008}. Each row $(\mathbf{H}_{\rm NMF})_k$ is a
non-negative optical-depth component; we negate and $\ell_2$-normalise to match the
PCA column convention,
\begin{equation}
  \phi_k^{\rm NMF} \;=\; -\,\frac{(\mathbf{H}_{\rm NMF})_k}{\|(\mathbf{H}_{\rm NMF})_k\|_2}
  \;\leq\; 0,
  \label{eq:nmfphi}
\end{equation}
producing non-positive templates that are zero in the continuum and negative in
the line cores, matching the convention of the PCA components exactly. The
resulting $(\mu, \{\phi_k^{\rm NMF}\})$ pair is a drop-in replacement for the PCA
basis in all downstream machinery.

\paragraph{Test-time inference with bounded least squares.}
The separable velocity scan and Newton refinement of
Sections~\ref{sec:fft}--\ref{sec:opt} proceed identically for NMF bases, because
the outer nonlinear scan and the inner linear solve are independent of how the
template columns were constructed. Non-negativity of the weights is enforced in a
\emph{post-refinement} step: once the nonlinear parameters $(v, R, p_{\rm tell})$
are fixed at their Newton optimum, the full linear weight vector $\boldsymbol\theta$
is re-solved under the bounded system
\begin{equation}
  \widehat{\boldsymbol\theta} \;=\; \argmin_{\boldsymbol\theta}\;
  \|\sqrt{\mathbf{C}^{-1}}\,(\mathbf{A}\,\boldsymbol\theta - \mathbf{Y})\|^2
  \quad\text{s.t.}\quad
  w_k \geq 0\;\;(k=0,\dots,K),\quad c_m \in \mathbb{R},
  \label{eq:bvls}
\end{equation}
solved by bounded-variable least squares (BVLS, \texttt{scipy.optimize.lsq\_linear})
with lower bounds $\ell_k=0$ on the template columns and $\ell_m=-\infty$ on the
continuum columns. The continuum coefficients remain unconstrained because the
Legendre polynomials carry no physical sign requirement. Any ridge penalty is
absorbed into the augmented system before the BVLS call. The post-step recomputes
$\chi^2$ so that the reported goodness-of-fit reflects the constrained solution.

\subsection*{Comparison}

\paragraph{PCA: optimal compression, unconstrained weights.}
For a given $K$, PCA minimises the reconstruction error over \emph{all}
$K$-dimensional subspaces, so it is the optimal linear compressor --- no other
$K$-component basis captures more variance from the training library. The
unconstrained WLS inner solve is exact and closed-form. However, because the
weights $w_k$ are free, the reconstructed rectified log-flux
$\hat{M}(x) = \sum_k w_k\phi_k(x)$ can be \emph{positive}: in linear flux this
corresponds to $e^{\hat{M}} > 1$, i.e.\ the model spectrum exceeds the continuum.
For a purely absorbing stellar atmosphere this is unphysical (emission is not
modelled), so the positivity excursions are an artefact of the unconstrained
weights acting on a non-orthogonal projection of the data.

\paragraph{NMF: physical bounds, sub-optimal compression.}
Because the NMF components $\phi_k^{\rm NMF}$ are non-positive and the weights are
constrained to $w_k \geq 0$, the reconstructed log-flux
$\sum_k w_k\phi_k^{\rm NMF} \leq 0$ everywhere, and the corresponding linear-flux
rectified spectrum lies in $(0,1]$: strictly between zero (total opacity) and the
continuum (no absorption). The model is therefore physically consistent by
construction. The cost is compression efficiency: the non-negativity constraint
restricts the feasible basis, so NMF requires more components than PCA to capture
the same fraction of variance, and the BVLS post-step is slower than the
closed-form WLS solve. The components are also no longer orthogonal
($\mathbf{M}_{TT}$ is generally denser), though this does not affect the
correctness of the separable framework.

% ============================================================================
\section{Radial velocity as translation: the per-shift linear solve}
\label{sec:rv}
% ============================================================================
A radial velocity $v$ Doppler-shifts every wavelength by the factor $1+v/c$
(non-relativistically). In log-wavelength this is a \emph{uniform translation},
\begin{equation}
  \lambda \to \lambda\,(1+v/c)
  \quad\Longleftrightarrow\quad
  x \to x + \dx,\qquad \dx \equiv \ln\!\Big(1+\frac{v}{c}\Big),
\end{equation}
which on the pixel grid is a shift of $p=\dx/\delta$ pixels. The model templates
therefore become $T_k(x)\to T_k(x-\dx)$, i.e.\ each column of $\mathbf{T}$ is
translated by $p$ pixels. We write the shifted design as $\mathbf{T}(p)$, with
columns $[T_k(x_i - p\,\delta)]_i$.

The forward model \eqref{eq:linmodel} now reads
$Y(x)=\sum_k w_k\,T_k(x-\dx)$. Crucially, it is \emph{linear in $\mathbf{w}$ only
at a fixed shift $p$}; the dependence on $p$ is nonlinear. This structure is
\emph{separable}: we scan the single nonlinear parameter $p$ on the
\emph{outside} and solve the linear weight problem exactly on the \emph{inside}.
For each trial shift,
\begin{equation}
  \mathbf{N}(p) = \mathbf{T}(p)^\T\mathbf{C}^{-1}\,\mathbf{T}(p),
  \qquad
  \mathbf{r}(p) = \mathbf{T}(p)^\T\mathbf{C}^{-1}\,\mathbf{Y},
  \qquad
  \widehat{\mathbf{w}}(p)=\mathbf{N}(p)^{-1}\mathbf{r}(p),
  \label{eq:shiftsolve}
\end{equation}
and the \emph{profiled} objective, obtained by minimising out $\mathbf{w}$ as in
\eqref{eq:chi2profiled}, is
\begin{equation}
  \chi^2(p)\;=\;\min_{\mathbf{w}}\chi^2(\mathbf{w},p)
  \;=\;\mathbf{Y}^\T\mathbf{C}^{-1}\mathbf{Y}-\mathbf{r}(p)^\T\mathbf{N}(p)^{-1}\mathbf{r}(p).
  \label{eq:profiledp}
\end{equation}
The radial-velocity estimate is the global minimiser,
\begin{equation}
  \widehat{p}\;=\;\argmin_{p}\,\chi^2(p),
  \qquad
  \widehat{v}\;=\;c\,\big(e^{\widehat{p}\,\delta}-1\big).
\end{equation}
For a single template ($K=0$, no continuum, uniform weights) the cross term
$\mathbf{r}(p)^\T\mathbf{N}(p)^{-1}\mathbf{r}(p)$ reduces to the squared
template--data cross-correlation, so maximising it over $p$ is exactly the
classical cross-correlation velocity measurement of \citet*{TonryDavis1979}.
Equation~\eqref{eq:profiledp} is its generalisation to a multi-column template
\emph{basis} (and, in Sections~\ref{sec:continuum}--\ref{sec:resolution}, to a
joint continuum and resolution model).

% ============================================================================
\section{FFT acceleration of the velocity scan}
\label{sec:fft}
% ============================================================================
A direct evaluation of \eqref{eq:profiledp} on a grid of $n_{\rm shifts}$ trial
velocities, each requiring $\mathbf{T}(p)^\T\mathbf{C}^{-1}\mathbf{T}(p)$ and
$\mathbf{T}(p)^\T\mathbf{C}^{-1}\mathbf{Y}$, costs
$O(n_{\rm shifts}\,n\,K^2)$. The expensive shift-dependent pieces are, however,
all \emph{cross-correlations of fixed vectors against a shifted template}, which a
single FFT per template column delivers at \emph{all} shifts simultaneously.

Consider the right-hand side. Its $k$-th template entry, as a function of the
integer lag $p$, is
\begin{equation}
  b_{T,k}(p) \;=\; \sum_{i} T_k(x_i - p\,\delta)\,C_i^{-1}\,Y_i
  \;=\; \big[\,T_k \star (\mathbf{C}^{-1}\mathbf{Y})\,\big](p),
  \label{eq:bTcorr}
\end{equation}
a discrete cross-correlation of the fixed weighted-data vector $\mathbf{C}^{-1}\mathbf{Y}$
against $T_k$. By the cross-correlation theorem, for any two real vectors
$t,s$ of length $n$,
\begin{equation}
  \operatorname{corr}(t,s)\;=\;\irfft\!\big(\,\overline{\rfft(t)}\,\cdot\,\rfft(s)\,\big),
  \label{eq:corrthm}
\end{equation}
where $\rfft$/$\irfft$ are the real forward/inverse FFTs, the product is
elementwise over frequencies $\omega$, and the overline denotes complex
conjugation. Hence
\begin{equation}
  \mathbf{b}_T(\cdot)\;=\;\irfft\!\Big(\,\overline{\rfft(\mathbf{T})}\,\cdot\,\rfft(\mathbf{C}^{-1}\mathbf{Y})\Big),
  \label{eq:bTfft}
\end{equation}
computed once for every lag and every template column. A fractional (sub-pixel)
shift is applied in the same domain as a linear phase ramp,
\begin{equation}
  \rfft\!\big[T_k(\cdot - p\,\delta)\big]
  \;=\; \rfft(T_k)\,\cdot\,e^{-2\pi i\,\omega\,p},
  \label{eq:phaseramp}
\end{equation}
so the model template at any real-valued $p$ is recovered by
$\irfft\!\big(\rfft(T_k)\,e^{-2\pi i\omega p}\big)$.

The normal-matrix block $\mathbf{M}_{TT}(p)=\mathbf{T}(p)^\T\mathbf{C}^{-1}\mathbf{T}(p)$
depends on $p$ \emph{only} through the (slowly varying) weights $\mathbf{C}^{-1}$: if
$\mathbf{C}^{-1}$ were constant, $\mathbf{M}_{TT}$ would be exactly shift-invariant.
We therefore approximate $\mathbf{M}_{TT}(p)\approx\mathbf{M}_{TT}(0)$ during the
scan (cheap, evaluated once) and recompute $\mathbf{M}_{TT}$ \emph{exactly} at the
refined optimum (Section~\ref{sec:opt}).

\paragraph{Circular-edge handling.}
The FFT correlation in \eqref{eq:corrthm} is \emph{circular}: content shifted off
one end wraps to the other. To make the circular correlation equal the desired
linear one over the region of interest, we zero the data weights $C_i^{-1}$ in a band
of width equal to the maximum lag $|p|_{\max}$ at \emph{both} ends of the grid.
Within the interior valid region, no wrapped pixel ever carries weight for any
evaluated shift, so the circular and linear correlations coincide.

\paragraph{Cost.}
The whole velocity scan now costs $O\!\big(K\,n\log n\big)$ --- a handful of FFTs
of length $n$, one per template column --- instead of
$O(n_{\rm shifts}\,n\,K^2)$, a decisive saving when $n_{\rm shifts}$ and $n$ are
both large.

% ============================================================================
\section{Joint RV and weights: the convexity structure}
\label{sec:convexity}
% ============================================================================
A natural question is whether the joint problem in $(\mathbf{w},p)$ is convex. The
answer is block-structured, and worth stating precisely because it dictates the
optimisation strategy.

\paragraph{Convex in the linear block.}
At any fixed shift $p$, the objective $\chi^2(\mathbf{w},p)$ in \eqref{eq:chi2lin}
(with $\mathbf{T}\!\to\!\mathbf{T}(p)$) is a convex quadratic in $\mathbf{w}$ with
Hessian $2\mathbf{N}(p)\succeq 0$; it is globally and (when $\mathbf{N}(p)\succ0$)
uniquely solvable in closed form by \eqref{eq:shiftsolve}. There are no spurious
local minima in the weights.

\paragraph{Non-convex in the nonlinear parameter.}
The profiled objective $\chi^2(p)$ of \eqref{eq:profiledp}, obtained by
minimising out $\mathbf{w}$, is \emph{not} convex in $p$. Up to the constant
$\mathbf{Y}^\T\mathbf{C}^{-1}\mathbf{Y}$, it is the (negated) generalised
cross-correlation surface $\mathbf{r}(p)^\T\mathbf{N}(p)^{-1}\mathbf{r}(p)$, which
is inherently \emph{multimodal}: the quasi-periodic absorption-line pattern and
the self-similarity of the templates produce sidelobes (secondary minima) at lags
offset from the truth by characteristic line spacings, with one dominant minimum
at the true velocity.

\paragraph{Consequence.}
The joint problem is thus \emph{convex in the linear block,
non-convex (multimodal) in the nonlinear parameter}. Globally it is not convex;
block-wise (in $\mathbf{w}$) and locally (near the dominant peak) it is. This is
exactly why the method uses a \emph{global coarse scan} over $p$ to identify the
basin of the dominant minimum (Section~\ref{sec:fft}), followed by a
\emph{local refinement} (Section~\ref{sec:opt}) in which the main peak is smooth
and approximately parabolic, hence locally well-behaved (locally convex). The
short answer to ``is it convex?'' is: \emph{no globally, yes block-wise and
locally}.

% ============================================================================
\section{Adding the continuum within the same linear model}
\label{sec:continuum}
% ============================================================================
The observed spectrum is the product of a smooth continuum and the rectified
(line) spectrum. In log-flux this product becomes a sum, so a smooth continuum is
an \emph{additive} term that we represent with low-order Legendre polynomials
$P_m$, $m=0,\dots,L$, evaluated on $x$ mapped affinely onto $[-1,1]$ across the
fit range,
\begin{equation}
  \tilde{x}_i \;=\; 2\,\frac{x_i - x_1}{x_n - x_1} - 1 \;\in\; [-1,1],
  \qquad
  (\mathbf{P})_{im}\;=\;P_m(\tilde{x}_i).
\end{equation}
The full forward model in log-flux is then
\begin{equation}
  Y(x) \;=\; \underbrace{\sum_{k=0}^{K} w_k\,T_k(x-\dx)}_{\text{rectified line basis}}
  \;+\; \underbrace{\sum_{m=0}^{L} c_m\,P_m(x)}_{\text{log-continuum}} .
  \label{eq:fullmodel}
\end{equation}
We augment the design matrix by appending the continuum block,
\begin{equation}
  \mathbf{A}(p) \;=\; \big[\,\mathbf{T}(p)\;\big|\;\mathbf{P}\,\big],
  \qquad
  \boldsymbol\theta \;=\; (\mathbf{w},\mathbf{c}),
\end{equation}
with $\mathbf{A}$ of size $n\times D$, $D=(K{+}1)+(L{+}1)$. The continuum block is
\emph{shift-independent}, so its contributions to the normal equations are
precomputed once,
\begin{equation}
  \mathbf{M}_{PP} = \mathbf{P}^\T\mathbf{C}^{-1}\mathbf{P},
  \qquad
  \mathbf{b}_P = \mathbf{P}^\T\mathbf{C}^{-1}\mathbf{Y}
  \qquad(\text{independent of }p).
\end{equation}
Only the template auto-block and the template--continuum cross-block depend on
$p$:
\begin{equation}
  \mathbf{M}_{TT}(p) = \mathbf{T}(p)^\T\mathbf{C}^{-1}\mathbf{T}(p),
  \qquad
  \mathbf{M}_{TP}(p) = \mathbf{T}(p)^\T\mathbf{C}^{-1}\mathbf{P},
\end{equation}
and the cross-block $\mathbf{M}_{TP}(p)$ is, like \eqref{eq:bTcorr}, a
cross-correlation of the fixed weighted-continuum columns
$\mathbf{C}^{-1}\mathbf{P}$ against the shifted templates --- so it too is obtained at
all lags by FFT, $\mathbf{M}_{TP}(\cdot)=\irfft\big(\overline{\rfft(\mathbf{T})}\cdot\rfft(\mathbf{C}^{-1}\mathbf{P})\big)$.
The block normal equations are
\begin{equation}
  \mathbf{N}(p) =
  \begin{bmatrix}
    \mathbf{M}_{TT}(p) & \mathbf{M}_{TP}(p)\\[2pt]
    \mathbf{M}_{TP}(p)^\T & \mathbf{M}_{PP}
  \end{bmatrix},
  \qquad
  \mathbf{r}(p) =
  \begin{bmatrix}
    \mathbf{b}_T(p)\\[2pt]
    \mathbf{b}_P
  \end{bmatrix},
  \label{eq:blockN}
\end{equation}
with the inner solve and profiled objective exactly as before,
\begin{equation}
  \widehat{\boldsymbol\theta}(p) = \mathbf{N}(p)^{-1}\mathbf{r}(p),
  \qquad
  \chi^2(p) = \mathbf{Y}^\T\mathbf{C}^{-1}\mathbf{Y} - \mathbf{r}(p)^\T\mathbf{N}(p)^{-1}\mathbf{r}(p).
  \label{eq:profiledp2}
\end{equation}

\paragraph{Why separate the continuum from the templates.}
Because the templates are \emph{rectified} (their pseudo-continuum has been
removed before the PCA), they carry essentially only line information and are
approximately zero in the continuum. The smooth flux level is carried entirely by
the orthogonal Legendre block. This separation of roles is what avoids the
classic \emph{template/continuum degeneracy}: were the templates allowed to carry
broadband shape, a low-order continuum could partially mimic and absorb it,
leaving $\mathbf{N}$ ill-conditioned and $(\mathbf{w},\mathbf{c})$ poorly
determined. Rectifying the templates and modelling the continuum as a separate,
smooth, additive block keeps the two contributions linearly independent (to the
extent that line structure and low-order polynomials are dissimilar) and the solve
well-conditioned. This is the same design principle as the additive and
multiplicative polynomial continua in pPXF \citep*{Cappellari2017} and the
PCA-template plus polynomial model of \citet*{Bolton2012}, here cast entirely in
log-flux so that the multiplicative continuum becomes additive.

% ============================================================================
\section{Adding the spectral resolution}
\label{sec:resolution}
% ============================================================================
A finite, lower instrumental resolution convolves the spectrum with a
line-spread function, which we take to be Gaussian. At constant resolving power
$R=\lambda/\Delta\lambda$, the LSF has a constant width \emph{in velocity},
$\mathrm{FWHM}_v = c/R$, hence a constant velocity dispersion
\begin{equation}
  \sigma_v \;=\; \frac{c}{f\,R},
  \qquad f \equiv 2\sqrt{2\ln 2}.
\end{equation}
On the log-wavelength grid, which is uniform in velocity with
$\velscale = c\,\delta$ km\,s$^{-1}$ per pixel, this is a constant width in
\emph{pixels},
\begin{equation}
  \sigma_{\rm pix} \;=\; \frac{\sigma_v}{\velscale} \;=\; \frac{c}{f\,R\,\velscale}.
  \label{eq:sigpix}
\end{equation}
A convolution by a Gaussian of fixed pixel width is a single multiply in the
Fourier domain: a Gaussian of standard deviation $\sigma_{\rm pix}$ has transform
$\exp(-2\pi^2\omega^2\sigma_{\rm pix}^2)$, so we broaden each template ``on the
fly'' by
\begin{equation}
  \widehat{T}_k(\omega) \;\longrightarrow\; \widehat{T}_k(\omega)\,
  \exp\!\big(-2\pi^2\,\omega^2\,\sigma_{\rm pix}^2\big).
  \label{eq:fftgauss}
\end{equation}
The Fourier shift \eqref{eq:phaseramp} and the Fourier Gaussian
\eqref{eq:fftgauss} compose trivially as a single elementwise product in the
frequency domain,
\begin{equation}
  \widehat{T}_k(\omega)\;\longrightarrow\;
  \widehat{T}_k(\omega)\,
  \exp\!\big(-2\pi^2\omega^2\sigma_{\rm pix}^2\big)\,
  e^{-2\pi i\omega p},
\end{equation}
so the broadened, shifted template at any real $(p,\sigma_{\rm pix})$ costs one
inverse FFT. In particular the $\mathbf{b}_T$ and $\mathbf{M}_{TP}$
cross-correlations at fixed $R$ are still FFTs --- broaden the template transform,
then correlate.

We promote $R$ (equivalently $\sigma_{\rm pix}$) to a \emph{second nonlinear,
outer} parameter alongside $p$. The structure is unchanged: at fixed
$(p,R)$ the model \eqref{eq:fullmodel} is still linear in
$\boldsymbol\theta=(\mathbf{w},\mathbf{c})$, so the same separable scheme applies
--- a $2$-D outer search/optimisation over $(p,R)$ wrapping the inner linear solve
\eqref{eq:profiledp2}.

\paragraph{An honest caveat: convolution and the logarithm do not commute.}
The LSF physically convolves the \emph{linear} flux. But the inner solve is linear
in $\boldsymbol\theta$ only if we broaden the \emph{log-rectified} templates,
which keeps the additive log-continuum intact. These are not the same operation,
because convolution does not commute with the logarithm. Writing the rectified
log-flux as $M(x)$ and the LSF as a normalised kernel, a cumulant (second-order)
expansion of the exact linear-flux broadening gives
\begin{equation}
  \ln\!\big(\mathrm{LSF}\!*\!e^{M}\big)
  \;\approx\; \mathrm{LSF}\!*\!M \;+\; \half\,\Var_{\mathrm{LSF}}(M)
  \;+\;\cdots,
  \label{eq:cumulant}
\end{equation}
where $\Var_{\mathrm{LSF}}(M)$ is the local variance of $M$ under the kernel. The
leading term $\mathrm{LSF}\!*\!M$ is precisely our log-space broadening
\eqref{eq:fftgauss}; the $\tfrac12\Var$ term is the first correction, which is
positive and largest where $M$ varies sharply --- i.e.\ in deep, saturated line
cores --- so log-space broadening slightly \emph{over-deepens} those cores at
second order. The continuum, being smooth on the scale of the LSF, has negligible
variance under the kernel and so factors out, $\mathrm{LSF}\!*\!(\text{smooth})
\approx\text{smooth}$, and remains additive --- exactly the property we need for
linearity. We therefore conclude that the log-space convolution is a
\emph{first-order-accurate} approximation: excellent for shallow and moderate
lines, with a small, bounded core bias for the deepest lines. The exact
linear-flux convolution, by contrast, would mix the components nonlinearly
(broadening $\sum_k w_k T_k$ in linear flux is not $\sum_k w_k$ times broadened
$T_k$ unless one broadens a \emph{single} fixed spectrum), breaking the linearity
of the inner solve for a multi-component basis.

\paragraph{Removing the core bias by iteration.}
The bias is, however, removable without giving up the linear inner solve, by a
fixed-point iteration. Write the log-space-broadened model block as $A_c =
\mathrm{LSF}\!*\!M_c$ and the exact one as $E_c = \ln(\mathrm{LSF}\!*\!e^{M_c})$,
where $M_c$ is the (un-broadened, shifted) log-flux of component $c\in\{\text{star,
tell}\}$ from the current solution. Their difference $\delta = \sum_c (E_c - A_c)$
is precisely the $\tfrac12\Var$-and-higher correction (concentrated in the cores).
Since the exact model is $E = A + \delta$, fitting the \emph{linear} machinery to
the corrected data $\mathbf{Y} - \delta$ makes the exact model match the data; as
$\delta$ depends on the solution we iterate $\delta\!\to\!\text{fit}\!\to\!\delta$
to a fixed point. Each pass is the usual linear solve (the normal matrix is
unchanged; only the right-hand side moves) plus one forward linear-flux
convolution, and the nonlinear parameters are re-optimised against the corrected
data, so the resolution bias is removed as well. On an in-basis test the reduced
$\chi^2$ falls from $\sim\!40$ to $\sim\!1$ and the recovered $R$ bias from
$\sim\!-2\%$ to $\lesssim 0.1\%$; the iteration converges geometrically (a handful
to $\sim\!10$ steps). This is the \texttt{n\_conv\_iter} option (off by default,
which leaves the fast first-order approximation untouched).

\paragraph{Beyond a single Gaussian.}
For a wavelength-dependent $R(\lambda)$ or a non-Gaussian LSF there is no
closed-form Fourier multiply. The practical remedy is to precompute banded
(sparse) convolution matrices on a grid of $R$ values and interpolate among them
at fit time, retaining the linear inner solve at the cost of the single-multiply
FFT shortcut.

% ============================================================================
\section{Tellurics and the instrument wavelength zero-point}
\label{sec:tellurics}
% ============================================================================
Telluric absorption from Earth's atmosphere is imprinted at \emph{fixed observed
(topocentric) wavelengths}, independent of the star's motion. We model it with a
second rectified basis $\mathbf{S}=[\,S_0\mid S_1\mid\cdots\mid S_J\,]$, built
exactly as the stellar templates (log-rectify $+$ PCA) but from a grid of telluric
transmission models. Because the tellurics do not share the stellar Doppler shift,
$\mathbf{S}$ is \emph{not} translated by $\dx$; like the continuum it enters the
design as shift-independent columns, broadened by the instrument LSF only (no
rotational broadening). The forward model \eqref{eq:fullmodel} becomes
\begin{equation}
  Y(x) \;=\; \underbrace{\sum_{k} w_k\,T_k(x-\dx)}_{\text{stellar}}
  \;+\; \underbrace{\sum_{j} u_j\,S_j(x)}_{\text{telluric (observed frame)}}
  \;+\; \underbrace{\sum_{m} c_m\,P_m(x)}_{\text{continuum}},
\end{equation}
with the design augmented to
$\mathbf{A}(p)=[\,\mathbf{T}(p)\mid\mathbf{S}\mid\mathbf{P}\,]$ and the linear vector
to $\boldsymbol\theta=(\mathbf{w},\mathbf{u},\mathbf{c})$. The telluric and continuum
blocks are shift-independent, so their normal-equation contributions are precomputed
once; only the stellar auto- and cross-blocks depend on $p$, and the separable
scan/solve is otherwise unchanged.

\paragraph{The wavelength zero-point.}
Real spectrographs carry a slit-dependent \emph{wavelength zero-point} error (e.g.\
uncorrected flexure on Keck/DEIMOS): an overall shift of the observed wavelength
scale by some lag $p_0$. The stellar template then matches the data at
$p_{\rm star}=p_{\rm true}+p_0$, conflating the true Doppler shift with the
instrumental zero-point. Because the tellurics are at rest in the observer frame,
their apparent shift measures $p_0$ \emph{directly}. We therefore optionally promote
a telluric lag $p_{\rm tell}$ to a further nonlinear parameter: the telluric columns
are translated by a Fourier phase ramp,
$\widehat{S}_j(\omega)\to\widehat{S}_j(\omega)\,e^{-2\pi i\omega p_{\rm tell}}$,
exactly as the stellar templates are translated by $p$ (Section~\ref{sec:fft}). At
fixed $(p,R,p_{\rm tell})$ the model is still linear in
$(\mathbf{w},\mathbf{u},\mathbf{c})$, so the same separable least squares applies.

Unlike the radial velocity $p$ --- whose profiled $\chi^2$ is multimodal and so
requires the global FFT scan of Section~\ref{sec:fft} --- the zero-point is small
and its surface is smooth and unimodal near $p_{\rm tell}=0$. We therefore do
\emph{not} add a coarse grid dimension for it: the coarse scan keeps the tellurics at
$p_{\rm tell}=0$ (in the shift-independent block) to locate the $(p,R)$ basin, and
$p_{\rm tell}$ is initialised at $0$ and refined \emph{jointly} with the other
nonlinear parameters by the autodiff Newton step of Section~\ref{sec:opt}.

\paragraph{Zero-point-corrected velocity.}
Since $p_{\rm star}=p_{\rm true}+p_0$ and $p_{\rm tell}=p_0$, the instrumental
zero-point cancels in the \emph{difference}, giving the corrected radial velocity
\begin{equation}
  \widehat v_{\rm corr} \;=\; c\,\Big(e^{(p_{\rm star}-p_{\rm tell})\,\delta}-1\Big),
  \label{eq:vcorr}
\end{equation}
with its uncertainty propagated from the joint $(p_{\rm star},p_{\rm tell})$ block of
the nonlinear-parameter covariance. A barycentric correction is then added to place
$\widehat v_{\rm corr}$ in the heliocentric frame. This reproduces the telluric/
sky-line wavelength recalibration used by DEIMOS pipelines (e.g.\ \texttt{dmost}).

% ============================================================================
\section{Optimisation and uncertainties}
\label{sec:opt}
% ============================================================================
The fit proceeds in two stages, matched to the convexity structure of
Section~\ref{sec:convexity}.

\paragraph{(1) Global coarse scan.}
The FFT machinery of Section~\ref{sec:fft} evaluates the profiled objective
\eqref{eq:profiledp2} on a dense grid of integer lags $p$ across the velocity
search window (and, when resolution is fit, on a coarse grid of $R$ values via the
broadening \eqref{eq:fftgauss}, using the cheap approximation
$\mathbf{M}_{TT}(p)\approx\mathbf{M}_{TT}(0)$). The minimum over this $1$-D (or
$2$-D) surface locates the basin of the dominant peak, immune to the sidelobes.

\paragraph{(2) Local refinement by autodiff Newton.}
Starting from the coarse minimiser, the profiled $\chi^2$ as a function of the
nonlinear parameters --- $p$, or $(p,\sigma_{\rm pix})$ when $R$ is fit --- is
minimised locally. Because the radial-velocity translation \eqref{eq:phaseramp},
the Gaussian broadening \eqref{eq:fftgauss}, and the inner linear solve
\eqref{eq:profiledp2} are all differentiable operations, the gradient and Hessian
of the profiled $\chi^2$ are obtained \emph{exactly} by JAX automatic
differentiation. The refinement is a damped (Levenberg--Marquardt) Newton
iteration: at each step the symmetrised Hessian is floored to positive definite by
LM damping, the resulting step is forced to be a descent direction, and its length
is set by a backtracking line search enforcing the Armijo sufficient-decrease
condition. On the smooth, low-dimensional, locally near-parabolic main peak this
converges in a few iterations.

\paragraph{Exact final solve.}
At the refined optimum $(\widehat p,\widehat\sigma_{\rm pix})$ the template
auto-block $\mathbf{M}_{TT}$ is recomputed \emph{exactly} (no $p=0$
approximation), the full design $\mathbf{A}(\widehat p)$ is assembled, and the
linear system \eqref{eq:profiledp2} is solved once more to deliver
$\widehat{\boldsymbol\theta}=(\widehat{\mathbf{w}},\widehat{\mathbf{c}})$, the
linear-parameter covariance $\mathbf{N}^{-1}$, and the exact $\chi^2$.

\paragraph{Uncertainties.}
The covariance of the nonlinear parameters follows from the curvature of the
profiled objective at the optimum. For a least-squares $\chi^2$ the standard
$\Delta\chi^2=1$ rule gives
\begin{equation}
  \operatorname{Cov}(\widehat p,\widehat\sigma_{\rm pix})
  \;=\; 2\,\big[\mathbf{H}\big]^{-1},
  \qquad
  \mathbf{H} = \nabla^2\chi^2\big(\widehat p,\widehat\sigma_{\rm pix}\big),
\end{equation}
the factor $2$ arising because $\tfrac12\mathbf{H}$ is the Fisher information of
the Gaussian likelihood $\propto\exp(-\chi^2/2)$. This covariance is propagated to
the physical parameters $(v,R)$ by the Jacobian of
$v=c(e^{p\delta}-1)$ and \eqref{eq:sigpix}; e.g.\ $\partial v/\partial p =
c\,\delta\,e^{p\delta}$ and $\partial R/\partial\sigma_{\rm pix} =
-R/\sigma_{\rm pix}$. When the data are consistent with no broadening
($\widehat\sigma_{\rm pix}\!\to\!0$), $R$ is reported as a one-sided lower limit
and only $v$ is refined.

\paragraph{Degrees of freedom.}
The reduced chi-square uses
\begin{equation}
  \mathrm{dof} \;=\; n_{\rm good} \;-\; (K+1) \;-\; (J+1) \;-\; (L+1) \;-\; n_{\rm nonlin},
\end{equation}
with $n_{\rm good}$ the number of unmasked pixels, $(K+1)$ template weights, the
$(J+1)$ telluric weights when a telluric basis is used (otherwise drop this term),
$(L+1)$ continuum coefficients, and $n_{\rm nonlin}$ the number of fitted nonlinear
parameters: the RV shift $p$, plus optionally the resolution $\sigma_{\rm pix}$, the
rotation $v\sin i$, and the telluric shift $p_{\rm tell}$.

\begin{thebibliography}{9}
\bibitem[Boutsidis \& Gallopoulos(2008)]{BoutsidisGallopoulos2008}
  Boutsidis, C., \& Gallopoulos, E.\ 2008,
  ``SVD based initialization: A head start for nonneg\-ative matrix factorization,''
  Pattern Recognition, 41, 1350.
\bibitem[Lee \& Seung(1999)]{LeeSeung1999}
  Lee, D.~D., \& Seung, H.~S.\ 1999,
  ``Learning the parts of objects by non-negative matrix factorization,''
  Nature, 401, 788.
\bibitem[Bolton et al.(2012)]{Bolton2012}
  Bolton, A.~S., Schlegel, D.~J., Aubourg, \'E., et al.\ 2012,
  ``Spectral Classification and Redshift Measurement for the SDSS-III
  Baryon Oscillation Spectroscopic Survey,''
  AJ, 144, 144.
\bibitem[Cappellari(2017)]{Cappellari2017}
  Cappellari, M.\ 2017,
  ``Improving the full spectrum fitting method: accurate convolution with
  Gauss--Hermite functions,''
  MNRAS, 466, 798.
\bibitem[Tonry \& Davis(1979)]{TonryDavis1979}
  Tonry, J., \& Davis, M.\ 1979,
  ``A survey of galaxy redshifts. I. Data reduction techniques,''
  AJ, 84, 1511.
\end{thebibliography}

\end{document}
